\section{Conclusion}\label{sec:conclusion}

We have proposed \emph{Topological Causal Data Analysis} as a
coherent research domain at the intersection of two mature
mathematical traditions. The case for the proposal rests on four
elements that the body of the paper has developed in turn: a formal
framework (\Cref{sec:framework}) anchored in the tuple
$(\Mc, G, \Tc, \Cc)$ of \Cref{def:tcda} and the foundational
propositions \Cref{prop:stability-transfer,prop:identifiability}; a
structured literature review (\Cref{sec:literature}) showing that the
methodological building blocks of TCDA are already in place but
fragmented across five non-interacting threads
(\Cref{sec:lit_gaps}); a suite of five case studies
(\Cref{sec:case_studies}) demonstrating that the framework is
non-vacuous and that each cell of its taxonomy admits an explicit
instantiation; and a research agenda (\Cref{sec:research}) comprising
eleven formal Open Questions, fourteen Research Avenues, and four
Computational Challenges calibrated against the gaps and case-study
limitations identified earlier.

\paragraph{The four cornerstones of the field.}
The argument of the paper is that the four seed works
$\{$\citet{ibeling2021topological}, \citet{kim2026topological},
\citet{lecci2014statistical},
\citet{elyaagoubi2023statistical}$\}$, taken together, span the
mathematical, statistical, and methodological corners of TCDA. The
first supplies the topology of causal-model spaces and the
no-free-lunch theorem that frames the field's negative results
\citep{ibeling2021topological}; the second supplies the
doubly-robust, stability-guaranteed estimator that converts a
topological summary into a calibrated causal quantity
\citep{kim2026topological}; the third supplies the statistical
inference scaffolding---confidence sets, bootstrap, rates of
convergence---without which the calibration of \emph{any} TCDA
estimator would be void \citep{lecci2014statistical}; and the fourth
supplies the simulation-based benchmark on which discovery procedures
in time-series and dependence-network settings can be evaluated
\citep{elyaagoubi2023statistical}. No single seed work is the field
on its own, and a serious TCDA programme must draw on all four
threads in parallel.

\paragraph{What we have not done.}
We have not claimed that the framework is final, that every
proposition is fully proved, that every open question is solvable on
its current statement, or that the taxonomy is mutually exclusive.
Several arguments in \Cref{sec:framework,sec:case_studies} are
proof sketches that demand fuller treatment, several research
avenues in \Cref{sec:research} are likely to fragment or merge as
the field matures, and the case studies of \Cref{sec:case_studies}
sit at varying levels of methodological readiness---some are
recapitulations of published work, others are concrete proposals
awaiting implementation. The four contributions enumerated in
\Cref{sec:intro_contributions} should therefore be understood as a
proposal in the technical sense---a unified vocabulary, a
taxonomy, a literature map, and an agenda---rather than as a
declaration of completion.

\paragraph{A call for interdisciplinary collaboration.}
The community most likely to settle the open questions of
\Cref{sec:research} does not yet exist. Its constituents---algebraic
topologists, statisticians, machine-learning researchers, causal-
inference theorists, computational scientists, and domain experts in
biomedicine, climate, finance, neuroscience, and materials---are
distributed across departments that historically have not shared
seminars or journals. The contribution that TCDA can make beyond its
internal technical agenda is to provide a vocabulary in which
collaborations across these silos can be staged: a topologist can
read \Cref{prop:stability-transfer} as a stability statement and
contribute new admissible functors $\Tc$; a causal-inference theorist
can read \Cref{prop:identifiability} as an identifiability question
and contribute new classes $\mathcal{G}$; a domain scientist can read
the case studies of \Cref{sec:case_studies} as templates and
contribute new applications; a computational scientist can read
\Cref{sec:research_compute} as a list of engineering targets. The
challenges that result---the GPU cubical PH of
\Cref{c:gpu-cubical}, the interoperable software of
\Cref{c:software}, the benchmark suite of \Cref{r:benchmarks}, the
validation standards of \Cref{r:validation-standards}---are
themselves invitations to a kind of work that no single sub-community
can complete in isolation. We hope this paper makes such invitations
easier to extend and easier to accept.

\paragraph{Closing.}
The slogan we offer is simple: \emph{the shape of data has a causal
story to tell}. Topology, on its own, recovers the shape; causal
inference, on its own, recovers the story; and only together can they
deliver, for the increasingly complex and increasingly non-Euclidean
data of contemporary science, the calibrated and intervention-aware
accounts that practical deployment demands. The four contributions of
this paper---a framework, a literature, a case-study suite, and an
agenda---are offered as the starting point of that joint enterprise,
not its conclusion. The body of work that converts the proposal into
a mature field remains to be done, and it is to that body of work
that we hope this paper now invites the reader.
